\section{Failure analysis}
\label{sec:formal-practice-05}

When a proposed proof fails, classify the failure before changing the theorem.
\begin{longtable}{@{}p{0.24\textwidth}p{0.66\textwidth}@{}}
\caption{Failure analysis}\\
\toprule
Failure & Meaning \\
\midrule
\endfirsthead
\caption[]{Failure analysis (continued)}\\
\toprule
Failure & Meaning \\
\midrule
\endhead
\midrule\multicolumn{2}{r}{\emph{continued on next page}}\\
\endfoot
\endlastfoot
No derivation found & The search may be incomplete; this alone gives no
semantic conclusion. \\
Semantic refutation constructed & Semantic consequence fails; by soundness,
derivability in the sound calculus also fails. \\
Rule side condition fails & That rule instance is unavailable; a different
derivation may exist. \\
Construction invariant fails & The proposed construction does not prove the
theorem; a corrected invariant or construction is required. \\
Hypothesis cannot be verified & The cited theorem cannot yet be applied. \\
Proof assistant rejects a term & The term does not type-check in the current
environment; the informal theorem may still be correct. \\
Solver returns unknown & No satisfiability or validity conclusion follows from
that response. \\
\bottomrule
\end{longtable}

\begin{exercise}
\label{ex:formal-practice-01}
Develop a countermodel to
\(\forall x(P(x)\lor Q(x))\entails
(\forall xP(x))\lor(\forall xQ(x))\).  Prove that a one-element domain cannot
serve and that a two-element domain can.
\end{exercise}

\begin{exercise}
\label{ex:formal-practice-02}
Write a proof plan for the substitution lemma.  Identify the induction object,
the strengthened induction statement, and the quantifier case in which bound
variable renaming is required.
\end{exercise}

